\section{Maximum-likelihood estimation}\label{sec:ml}

\subsection{Estimator}

Stack the real and imaginary parts of all 18 complex measurements of a voxel into
$y\in\mathbb R^{36}$ and the model \eqref{eq:signal} into $s(\theta)$. Under the Gaussian noise
model the negative log-likelihood is $\|r(\theta)\|^2/(2\sigma^2)$ with $r=s(\theta)-y$, so the ML
estimate is the nonlinear least-squares solution $\hat\theta=\arg\min_\theta\|r(\theta)\|^2$,
independent of $\sigma$. The model is the exact steady state \eqref{eq:iso}, so the estimator
and the CRLB use the same forward model.

\paragraph{Constraints.} The fit is a projected Levenberg--Marquardt (LM) iteration. After
every step the parameters are projected onto the physical domain of the Lorentzian model,
$\log\Ttwos\le\log\Ttwo$ ($\Rtwop\ge0$), and onto a box $\Bone\in[\E^{-1},540/330]$,
$\Tone\in[10,20\,000]$\,ms, $\Ttwo\in[1,5000]$\,ms. The upper $\Bone$ bound is the assumption
$\alpha_2\le540^\circ$ of~\cite{cheng2019}; it excludes the exact complex alias of
\cref{prop:equiv} by construction. Active constraints are recorded per voxel; such estimates
are not locally unbiased and fall outside the interior CRLB.

\paragraph{Initialisation.} Magnitude parameters and $\dw$ come from a sequential estimator;
the global phase has a closed-form optimum for fixed remaining parameters,
$\phi_0^{(0)}=\arg\bigl(\sum S_{ijk}\overline{S^{(0)}_{ijk}}\bigr)$, where $S^{(0)}$ is the model at the
initial values with $\phi_0=0$.

\paragraph{Iteration.} With $J=\partial r/\partial\theta$ (exact, \cref{sec:jacobian}),
\begin{equation}
  \bigl(J^{\mathsf T}J+\lambda\,\mathrm{diag}(J^{\mathsf T}J)\bigr)\Delta=-J^{\mathsf T}r,\qquad
  \theta\leftarrow\Pi(\theta+\Delta),
\end{equation}
where $\Pi$ is the projection. A step is accepted if it lowers $\|r\|^2$ ($\lambda\leftarrow0.3\lambda$),
otherwise rejected ($\lambda\leftarrow10\lambda$); $\lambda_0=10^{-2}$, $\lambda\in[10^{-9},10^9]$, 40 iterations.
Marquardt's diagonal scaling makes the step invariant to rescaling individual parameters. A
parameter held fixed (e.g.\ a given $\Bone$ map) has its Jacobian column removed and a unit
diagonal, so that $\Delta=0$ for it. For zero-residual (noise-free) problems with a nonsingular
Jacobian, Gauss--Newton converges locally quadratically, and noise-free fits recover the
parameters to $10^{-15}$; with noise the residual at the solution is nonzero and the local
convergence is in general linear, with a rate set by the residual and the model curvature.

\paragraph{Multiple starts.} Nonlinear least squares can stop at a local minimum: across the
$360^\circ$ branch, at the branch edge where scan 2 vanishes, or along the weak direction of
\cref{fig:eigen}. We run four starts---the model-fit estimate, the analytic inverse where
valid, and the model-fit estimate with $\Tone$ replaced by 700 and by 2000\,ms---and keep the
lowest final cost. In simulation the global optimum can be checked: a final cost above the
cost at the true parameters is certainly not the ML estimate.

\paragraph{Is one start enough?} On 2000 voxels with random tissue
($\Tone$ 600--3000\,ms, $\Ttwo$ 40--140\,ms, $\Ttwos/\Ttwo$ 0.6--0.95, $\Bone$ 0.85--1.15), a single start
ends at a higher cost than the best of four in \StartWorseThreeHundred\% of voxels at
SNR$(\Mzero)=300$ (median $|\Delta\log\hat\Tone|=\StartMedianThreeHundred$, maximum \StartMaxThreeHundred)
and \StartWorseThousand\% at SNR 1000. A $\chi^2$ test on the final cost (99.9\,\% quantile,
29 degrees of freedom) flags only \StartFlaggedThreeHundred\% and \StartFlaggedThousand\% of
these, respectively. One start with a cost check is therefore not sufficient; we use four
starts throughout, which the exact Jacobian makes affordable.

\subsection{Efficiency}

\Cref{tab:mc,tab:mcb1,tab:mcdiag} compare complex ML with the analytic inverse, a sequential
model-fit baseline (the same steps with $\Bone$, $\Tone$, $\Mzero$ fitted to the FID and echo
amplitudes of both scans) and the CRLB, on 1000 noise realisations of the reference tissue.
All ratios use ordinary standard deviations; the Monte Carlo standard error of an SD from
1000 samples is about $2\%$. (These runs used the box $\Bone\le\E$; for the published protocol no
estimate reached it, so the results are the same with the $540/330$ cap.)

\begin{table}[h]
\centering
\caption{Monte Carlo robust standard deviation (\%) of the log-estimates, 1000 noise realisations, reference tissue, $\Bone=1$. SNR is relative to $\Mzero$ (peak image SNR $\approx0.088\times$). Last column: fraction of voxels with undefined $\Tone$.}
\label{tab:mc}
\small
\begin{tabular}{lllccccc}
\toprule
Protocol & SNR$(\Mzero)$ & Estimator & $\Tone$ & $\Bone$ & $\Ttwo$ & $\Mzero$ & $\Tone$ undefined\\
\midrule
$15^\circ/330^\circ$ & 300 & CRLB & 13.5 & 1.0 & 7.7 & 6.3 & --\\
 &  & Joint ML & 13.7 & 0.9 & 7.5 & 6.0 & 0.0\%\\
 &  & Model fit & 90.7 & 3.7 & 36.8 & 56.9 & 0.0\%\\
 &  & Analytic inverse~\cite{cheng2019} & 102.6 & 46.1 & 37.0 & 57.1 & 1.9\%\\
 &  & Analytic inverse, $\Bone$ given & 35.9 & -- & -- & 23.3 & 1.9\%\\
\midrule
$15^\circ/330^\circ$ & 1000 & CRLB & 4.1 & 0.3 & 2.3 & 1.9 & --\\
 &  & Joint ML & 3.9 & 0.3 & 2.5 & 1.9 & 0.0\%\\
 &  & Model fit & 11.4 & 1.3 & 11.0 & 8.5 & 0.0\%\\
 &  & Analytic inverse~\cite{cheng2019} & 12.7 & 1.5 & 11.2 & 7.6 & 0.0\%\\
 &  & Analytic inverse, $\Bone$ given & 10.7 & -- & -- & 6.5 & 0.0\%\\
\midrule
$15^\circ/330^\circ$ & 3000 & CRLB & 1.4 & 0.1 & 0.8 & 0.6 & --\\
 &  & Joint ML & 1.3 & 0.1 & 0.8 & 0.6 & 0.0\%\\
 &  & Model fit & 4.0 & 0.4 & 3.8 & 2.9 & 0.0\%\\
 &  & Analytic inverse~\cite{cheng2019} & 4.4 & 0.4 & 3.7 & 2.6 & 0.0\%\\
 &  & Analytic inverse, $\Bone$ given & 3.9 & -- & -- & 2.2 & 0.0\%\\
\midrule
$15^\circ/30^\circ$ & 1000 & CRLB & 210.6 & 102.3 & 2.3 & 104.8 & --\\
 &  & Joint ML & 211.7 & 103.7 & 2.2 & 105.3 & 0.0\%\\
 &  & Model fit & 287.6 & 138.0 & 11.0 & 149.0 & 0.0\%\\
\midrule
$15^\circ/30^\circ$ & 3000 & CRLB & 70.2 & 34.1 & 0.8 & 34.9 & --\\
 &  & Joint ML & 77.1 & 37.3 & 0.7 & 38.1 & 0.0\%\\
 &  & Model fit & 225.6 & 109.7 & 3.8 & 113.9 & 0.0\%\\
\bottomrule
\end{tabular}
\end{table}


\begin{itemize}
\item \textbf{Where its assumptions hold, ML attains the bound.} At SNR 1000 and 3000 the
  SD/CRLB ratio is 0.99--1.02 for $\Tone$, $\Bone$, $\Ttwo$ and $\Mzero$, the bias is below $10^{-3}$ in
  log units, at most $0.3\%$ of estimates touch a constraint, and no voxel fails the cost
  check. The published protocol can separate $\Bone$ and $\Tone$.
\item \textbf{At SNR 300 the comparison is approximate.} The ratio is 1.00 for $\Bone$ and 0.95
  for $\Tone$, but $\Rtwop\ge0$ is active in $19.5\%$ of voxels (true $\Rtwop$ is small); those estimates
  are biased, which is how the $\Tone$ SD can fall slightly below the bound.
\item \textbf{The analytic inverse is 3.2--7.1 times above the bound} for $\Tone$, with bias
  $+0.44$ in $\log\Tone$ and $1.9\%$ undefined estimates at SNR 300; $\Bone$ is 4.8--28 times above.
  Given the true $\Bone$, its $\Tone$ step is 1.9--2.8 times above its matched bound (scan 1,
  $\Bone$ known: 57.1). The model-fit baseline, which shares the sequential structure but not
  the closed forms, behaves similarly, so the loss comes from the sequential structure.
\item \textbf{The small-angle control supports no efficiency claim.} For $15^\circ/30^\circ$,
  $10$--$35\%$ of ML estimates are on a box bound and biased; the bound itself is of order one
  in log units ($2.1$ at SNR 1000, a one-sigma factor of about 8), while $\Ttwo$ is unaffected.
\end{itemize}
